
\firstsection{Introduction}

\maketitle

Composite visualizations combine multiple charts to compare data, reveal
cross-view relationships, and communicate multivariate patterns. They are
common in visual analytics and communication~\cite{vasurvey, Survey16,
deng2023survey, AI4VIS}, including concatenation, layering, faceting, and
nesting~\cite{javed2012exploring, Composite, vaid}. Authoring them requires
coordinating data, encodings, layouts, coordinate systems, and dependencies
across charts.

Existing authoring approaches make trade-offs between predictability
and flexibility. Low-level libraries~\cite{D3} and graphical editors~\cite{Lyra,DataIllustrator,Charticulator} provide
fine-grained control, but require substantial effort to construct and
coordinate visual structures.
Explicit visual templates~\cite{ManyEyes,
RAWGraphs} instead expose recognizable chart forms with bounded
data roles, making intended outcomes easy to anticipate. Grammar- and schema-based approaches~\cite{VegaLite,Tableau,ApacheECharts} increase expressiveness by
abstracting across chart classes through shared marks, encoding channels, and
structural roles. However, this
flexibility requires authors to reason about increasingly abstract visual
structures, while explicit templates remain difficult to scale to the diverse
structures of composite visualizations.

We retain the clarity of explicit templates, but decompose composite
visualizations into predictable and recomposable chart blocks. Each block has
a recognizable visual structure and bounded behaviors, making the authoring
unit and expected outcome explicit. Related blocks are further organized
into chart families, providing controlled transitions across compatible
structures without abandoning their visual identities. The key question is
where to place these block and family boundaries: units that are too coarse
remain restrictive, whereas overly fine units lose the efficiency and
predictability of template-based authoring. We therefore conducted a formative
study in which visualization experts decomposed composite visualizations into
constituent units, informing the boundaries and organization of our chart
blocks and families.


Given these explicit units, the next question is how authors should express
composition relationships among them. Prior work characterizes composite
visualizations through recurring spatial relationships~\cite{
javed2012exploring,Composite,vaid}. We use these relationships directly as
references for interaction: dragging one block relative to another selects
layout-referenced targets that express concatenation, layering, faceting, or
nesting. Each action creates an explicit and editable composition relationship,
rather than merely determining the placement of a chart.

Explicit chart blocks nevertheless impose bounded data roles, which can make
exploration brittle when dimensions are added, removed, or reorganized. This
is particularly relevant to composite visualization, where a dimension may be
represented within a chart, distributed across repeated views, or expressed
through a composition relationship.

We therefore allow dimensional changes to move across chart-level and
composition-level representations. Related blocks within a chart family and
compatible composition structures can provide alternative realizations of the
same dimensional organization, while preserving compatible data assignments
and visual intent. This extends the flexibility of explicit chart blocks
without requiring authors to reconstruct their designs through lower-level
graphical abstractions.



Guided by these principles, we developed {\system}, an interactive authoring
tool for constructing composite visualizations from reusable chart blocks.
{\system} combines visually explicit chart units, layout-referenced
drag-and-drop composition, and data-aware dimensional adaptation to support
concatenation, layering, faceting, and nesting while keeping constituent
blocks, their relationships, and dimensional adaptations editable
throughout authoring.
We evaluate {\system} through case studies, a comparative user study, 
and an exploratory
user study spanning a broader range of composite visualization scenarios.
Together, these evaluations characterize how {\system} support a
consistent and learnable workflow for composite visualization authoring.


In summary, our contributions are:
\begin{itemize}
    \item A \textbf{formative study} that grounds the decomposition and granularity
    of reusable chart blocks in how visualization experts perceive composite
    visualizations;


    \item An \textbf{interactive authoring approach} that supports
     drag-and-drop composition with a data-aware engine
    for dimensional adaptation;


    \item A \textbf{comprehensive evaluation} of {\system}, including case studies, a comparative user study, and an exploratory user study.
\end{itemize}


